\documentclass[aspectratio=169,12pt]{beamer}

\usepackage{fontspec}
\usepackage{polyglossia}
\setdefaultlanguage{russian}
\setotherlanguage{english}

\setmainfont{Times New Roman}
\setsansfont{Arial}

\usepackage{ragged2e}
\usepackage{hyperref}
\hypersetup{hidelinks}

\setbeamertemplate{navigation symbols}{}
\setbeamersize{text margin left=8mm,text margin right=8mm}

% Читаемый номер и название на каждом рабочем слайде.
\setbeamerfont{footline}{size=\small}
\setbeamercolor{footline}{fg=black,bg=white}
\setbeamertemplate{footline}{%
  \leavevmode
  \hbox{%
    \begin{beamercolorbox}[wd=.72\paperwidth,ht=3.0ex,dp=1.2ex,leftskip=1em]{footline}
      \insertshorttitle{} --- \insertsectionhead
    \end{beamercolorbox}%
    \begin{beamercolorbox}[wd=.28\paperwidth,ht=3.0ex,dp=1.2ex,rightskip=1em plus1fil]{footline}
      Слайд \insertframenumber
    \end{beamercolorbox}%
  }%
}

\title[Согласование темы ВКР]{Согласование темы ВКР}
\subtitle{Автоматизация проверки документов с ИИ-агентами}
\author{Студент: Дениченко Александр}
\author{Научный сотрудник: Лисицина Василиса Васильевна}
\institute{Университет ИТМО}
\date{\today}

\begin{document}

% Слайд 0: тему можно оставить без номера.
\begin{frame}[plain,noframenumbering]
  \titlepage
\end{frame}

\section{Проблема и актуальность}
\begin{frame}{Слайд 1. Проблема и актуальность}
\justifying
\textbf{Текущая проблема:} ручной нормоконтроль проектных документов медленный, трудоемкий и зависит от опыта конкретного сотрудника.
\vspace{0.6em}
\begin{itemize}
  \item требования и артефакты проверки хранятся децентрализованно и в разных форматах;
  \item при параллельной работе по нескольким проектам резко растет число пропусков и несогласованностей;
  \item высокий порог входа для новых сотрудников, длительная адаптация;
  \item отсутствует единый, воспроизводимый и масштабируемый цикл проверки.
\end{itemize}
\vspace{0.6em}
\textbf{Почему это важно для ВКР:} автоматизация цикла проверки может одновременно повысить качество, скорость и управляемость процесса.
\end{frame}

\section{Цель и задачи}
\begin{frame}{Слайд 2. Цель и задачи исследования}
\textbf{Цель:} автоматизировать проверку документов на соответствие нормам с помощью ИИ-агентов, чтобы ускорить цикл контроля и повысить качество результата.

\textbf{Почему это нужно:}
\begin{itemize}
  \item сокращение времени проверки и нагрузки на экспертов;
  \item снижение зависимости результата от человеческого фактора.
\end{itemize}

\textbf{Ключевые задачи:}
\begin{itemize}
  \item анализ ролей и этапов проверки (интервью с экспертами);
  \item выбор инструментов: docx-парсинг, лингвистика, RAG и векторная БД;
  \item реализация сервиса с версионированием и процессом согласования;
  \item оценка эффектов: ускорение, экономия, порог входа, окупаемость.
\end{itemize}
\end{frame}

\section{Архитектура решения}
\begin{frame}{Слайд 3. Концепция и архитектура решения}
\begin{itemize}
  \item \textbf{Подход:} гибридный нормоконтроль (лингвистика + семантика + эксперт в контуре Human-in-the-Loop).
  \item \textbf{Ядро:} RAG-сервис с интеграцией LLM для проверки соответствия требованиям.
  \item \textbf{Компоненты:}
  \begin{itemize}
    \item docx-парсер и модуль преобразования в структурированное представление;
    \item сервис лингвистических проверок (орфография/пунктуация);
    \item векторизация требований и документов + векторная БД (ChromaDB);
    \item версионирование документов и хранение артефактов (MinIO);
    \item оркестрация процесса согласования (BPMN/Camunda).
  \end{itemize}
  \item \textbf{Результат:} воспроизводимый цикл проверки с объяснимыми замечаниями и журналом изменений.
\end{itemize}
\end{frame}

\section{Прогресс}
\begin{frame}{Слайд 4. Текущий прогресс по ВКР}
\textbf{Сделано на текущем этапе:}
\begin{itemize}
  \item сформулирована и детализирована тема ВКР;
  \item описаны бизнес-контекст, роли и основные сценарии проверки;
  \item подготовлены функциональные/нефункциональные требования;
  \item определен предварительный стек (RAG, LLM, ChromaDB, MinIO, BPMN).
\end{itemize}

\vspace{0.2em}
\textbf{Ближайшие шаги:}
\begin{itemize}
  \item зафиксировать baseline ручной проверки на интервью с экспертами;
  \item формализовать архитектуру взаимодействия агента с базой знаний;
  \item перейти к прототипу и экспериментам по времени/точности.
\end{itemize}
\end{frame}

\section{Ожидаемые результаты}
\begin{frame}{Слайд 5. Планируемые результаты и метрики}
\begin{itemize}
  \item прототип системы автоматизированной проверки проектных документов;
  \item методика сравнения ручной и автоматизированной проверки;
  \item количественные метрики:
  \begin{itemize}
    \item время проверки одного документа;
    \item полнота и точность обнаружения нарушений;
    \item сокращение трудозатрат и снижение порога входа;
    \item годовая экономия, эксплуатационные расходы, срок окупаемости.
  \end{itemize}
  \item практический артефакт: конфигурируемый процесс согласования с журналом версий.
\end{itemize}
\end{frame}

\section{Организация прослушивания}
\begin{frame}{Слайд 6. Формат прослушивания и запрос к комиссии}
\textbf{Цель текущего выступления:} согласование (или переформулирование) темы ВКР для приказа.

\vspace{0.6em}
\textbf{Организационно:}
\begin{itemize}
  \item доклад до 3 минут, далее вопросы/рекомендации комиссии;
  \item выступление, вероятно, в Kontur (не Zoom): заранее проверить демонстрацию экрана и звук;
  \item камера на этом этапе не требуется;
  \item будет вестись запись.
\end{itemize}

\vspace{0.6em}
\textbf{Прошу рекомендации комиссии:}
\begin{itemize}
  \item по финальной формулировке темы;
  \item по приоритету метрик эффективности;
  \item по объему пилота и границам апробации в ВКР.
\end{itemize}
\end{frame}

\end{document}

